\documentclass{article}
\begin{document}
\section{Das Grundgesetz} 
Das deutsche Grundgesetz wurde unter den Vorgaben der westlichen Besatzungsmächte vom Parlamentarischen Rat von 61 Politikern und 4 Politikerinnen in Blick geschrieben um aus den Fehlern der Verfassung der Weimarere Republik zu lernen. Zuerst war es nur als Übergangslösung gedacht, bis sie nach der Wiedervereinigung durch eine echte Verfassung ersetzt wird, dies ist aber nie passiert.

Im Grundgesetz, in mehrere Abschnitte aufgeteilt, werden unter anderem die Grundrechte der BürgerInnen beschrieben, die Länder beschrieben und die Staatsorgane begründet.

\subsection{Die Niedersächsische Verfassung}
In 1993 wurde die Niedersächsische Verfassung (Verf ND) verabschieded um die vorläufige Version auf 1951 abzulösen.

\end{document}